\PassOptionsToPackage{dvipsnames,svgnames,table}{xcolor}
\documentclass[10pt]{article}
\usepackage[a4paper,margin=22mm]{geometry}
\usepackage[no-math]{fontspec}\setmainfont[SmallCapsFont={Latin Modern Roman Caps}]{Latin Modern Roman}
\usepackage{amsmath,amssymb,amsthm,mathtools,graphicx,xcolor,hyperref,xurl,xspace,ifthen,enumerate,textcomp}
\usepackage[shortlabels]{enumitem}
\setlength{\emergencystretch}{4em}
\newtheorem{theorem}{Theorem}[section]
\newtheorem{theo}[theorem]{Theorem}
\newtheorem{thm}[theorem]{Theorem}
\newtheorem{lemma}[theorem]{Lemma}
\newtheorem{lemm}[theorem]{Lemma}
\newtheorem{prop}[theorem]{Proposition}
\newtheorem{coro}[theorem]{Corollary}
\newtheorem{corollary}[theorem]{Corollary}
\newtheorem{fact}[theorem]{Fact}
\newtheorem{claim}[theorem]{Claim}
\theoremstyle{definition}
\newtheorem{defi}[theorem]{Definition}
\newtheorem{definition}[theorem]{Definition}
\newtheorem{rema}[theorem]{Remark}
\newtheorem{remark}[theorem]{Remark}
\newtheorem{exam}[theorem]{Example}
\newtheorem{example}[theorem]{Example}
\newenvironment{enonce*}[2][]{\par\medskip\noindent\textbf{#2.}\itshape}{\par\medskip}
\providecommand{\eg}{e.g.\ }
\providecommand{\ie}{i.e.\ }
\providecommand{\resp}{resp.\ }
\providecommand{\cf}{cf.\ }
\providecommand{\longthanks}[1]{\par\smallskip\noindent #1\par\smallskip}
\newtheorem{proposition}[theorem]{Proposition}
\newtheorem{conjecture}[theorem]{Conjecture}
\numberwithin{equation}{section}
\newcommand{\R}{{\mathbb R}}
\newcommand{\Z}{{\mathbb Z}}
\newcommand{\N}{{\mathbb N}}
\newcommand{\expref}[2]{\texorpdfstring{\hyperref[#2]{#1~\ref{#2}}}{#1~\ref{#2}}}
\newcommand{\epfmtdoi}[1]{\href{https://doi.org/#1}{doi:#1}}
\newcommand{\epfmt}[2]{#1:\,#2}
\def\C{\mathbb{C}}
\def\E{\mathbb{E}}
\def\F{\mathbb{F}}
\def\eps{\varepsilon}
\def\R{\mathbb{R}}
\def\Var{\mathrm{Var}}
\newcommand{\ip}[1]{\left<#1\right>}
\def\P{\mathcal{P}}
\def\GL{\textrm{GL}}
\def\one{\mathbf{1}}
\def\Fr{\textrm{Fr}}
\def\sd{\textrm{dist}}

\newcommand{\cclass}[1]{\textsf{#1}}
\newcommand{\restate}[2]{\medskip \noindent {\bf #1} {\sl #2}}
\begin{document}
\begingroup\raggedright
\section*{1558. Almost-uniform product group actions support exact uniformity on the same support}
Noga Alon and Shachar Lovett. \emph{Almost k-Wise vs. k-Wise Independent Permutations, and Uniformity for General Group Actions}. Theory of Computing volume 9 article 15 (2013), official complete journal PDF acquired 2026-09-30; canonical DOI 10.4086/toc.2013.v009a015 verified against official landing and printed PDF header. \url{https://theoryofcomputing.org/articles/v009a015/}. CC BY3.0.
\paragraph{Selected problem and scope (editorial).} Exact original Theorem1.6 only: a distribution supported on S in a finite group G acting on finite X, almost (X times X)-uniform with error epsilon<0.5|X|\textasciicircum{}\{-7\}, supports an exactly X-uniform distribution on the same S. Complete original group-action definitions/statistical-distance convention, irreducible representation uniformity correspondence, convex-hull support equivalence, separator projection, coefficient/supremum bounds, product-representation moment control and full contradiction closing are retained. Named Maschke, Schur, convex separation and elementary finite-group representation/linear-algebra facts are prerequisites. The distribution-correction target1557, random-set existence and permutation parameter specializations are unselected and uncounted.
\paragraph{Difficulty (editorial).} After named Maschke/Schur and convex-separation inputs, the target must still project a support-separating functional onto the nontrivial irreducible coordinates, normalize its coefficient and supremum bounds, and transfer both moments through the X-times-X tensor representation with explicit error control. The complete authored separator/tensor-moment argument is native975–1193; it establishes exact feasibility on the unchanged support, which the nearby small-distance correction construction does not provide. These finite-dimensional support and quantitative moment obligations form one H2 task beyond routine Schur orthogonality or a direct convex-hull criterion.
\paragraph{Formatting (editorial).} Exact original human mathematical body and order retained; standalone typography, counter restoration and declared noncore printed locators only. Original comment prose excluded with percent guards retained; source typography and counters restored without upstream class execution. No mathematical repair or reconstruction.
\paragraph{Original source quirks (editorial).} The source explicitly defines statistical/total-variation distance as the unhalved l1 sum at native511; this convention is retained. The exact closing proof starts with epsilon<=0.5|X|\textasciicircum{}\{-7\} at1174 but closes under the original strict epsilon< threshold at1192, matching the selected theorem. Native1117 retains its original product wording while the adjoining tensor display includes complex conjugation; original wording and formulas are left unchanged without mathematical refereeing.
\endgroup
\clearpage
\setcounter{section}{1}\setcounter{theorem}{5}
\begin{theorem}\label{thm:almost_support_uniform}
Let $\mu$ be a distribution supported on a set $S \subset G$ which is
almost $(X \times X)$-uniform with error $\eps < 0.5 |X|^{-7}$. Then
there exists a distribution $\mu'$ supported on $S$ which is $X$-uniform.
\end{theorem}
\setcounter{section}{1}
\section{Preliminaries}
\label{sec:prelim}

For an integer $t \ge 1$ we denote $[t]:=\{1,2,\ldots,t\}$. The statistical (or total variation) distance between two distributions $\mu,\mu'$ is
given by $\sd(\mu,\mu')=\sum_{x} |\mu(x)-\mu'(x)|$. The expression $\delta_{i,j}$ is equal to $1$ if $i=j$ and is equal to $0$ otherwise.



\paragraph{Group action and uniformity}
A group $G$ \emph{acts} on a set $X$ if there is a homomorphism from $G$ to the
permutation group on $X$. That is, each $g \in G$ is a permutation on $X$,
and $(gh)(x)=g(h(x))$ for all $g,h \in G, x \in X$. We denote by $U_G$
the uniform distribution over $G$. For a distribution $\mu$ over $G$ we denote
by $g \sim \mu$ a random element chosen according to $\mu$. We abbreviate $\Pr_{g \in G}[\cdot] = \Pr_{g \sim U_G}[\cdot]$.

We recall some definitions from the
introduction: a distribution $\mu$ over $G$ is \emph{$X$-uniform} if
$$
\Pr_{g \sim \mu}[g(x)=y] = \Pr_{g \in G}[g(x)=y]
$$
for all $x,y \in X$; and a distribution $\mu$ is
\emph{almost $X$-uniform with error $\eps$} if
$$
\left| \Pr_{g \sim \mu}[g(x)=y] -
\Pr_{g \in G}[g(x)=y] \right| \le \eps
$$
for all $x,y \in X$. A family $S \subset G$ is \emph{$X$-uniform} (\emph{almost
$X$-uniform}, accordingly) if the uniform distribution over $S$ is such. We will need
some basic facts in linear algebra, geometry and representation theory,
which are presented below.

\paragraph{Linear algebra}
Let $A=(a_{i,j})$ be a complex matrix. The $L_{\infty}$ norm of $A$ is $\|A\|_{\infty}=\max
|a_{i,j}|$ (not to be confused with the operator norm of $A$ on $L_{\infty}$, which is not used in this paper).
The Frobenius norm of $A$ is $\|A\|_{\Fr} = \sqrt{\sum
|a_{i,j}|^2}$. For every $A$, $\|A\|_{\infty} \le \|A\|_{\Fr}$. A matrix
$A$ is unitary if $A \overline{A^t} = I$. The Frobenius
norm of a matrix is invariant under any unitary
basis change. That is, if $U$
is unitary then $\|U^{-1} A U\|_{\Fr} = \|A\|_{\Fr}$. The tensor product
of a $d_1 \times d_1$ matrix $A_1$ and a $d_2 \times d_2$ matrix $A_2$,
denoted $A_1 \otimes A_2$, is a $(d_1 d_2) \times (d_1 d_2)$ matrix, whose
entries are given by $(A_1 \otimes A_2)_{(i,i'),(j,j')} = (A_1)_{i,j}
(A_2)_{i',j'}$. The tensor product of unitary matrices is
also unitary.

\paragraph{Geometry}
Let $X=\{x_1,\ldots,x_N\}$ be a set of points in $\R^d$. The \emph{convex hull} of $X$ is the set of points contained in the minimal convex
set containing $X$; equivalently, it is the set of all points 
\[
\Bigl\{\sum
\lambda_i x_i: \lambda_1,\ldots,\lambda_N \ge 0, \;\sum \lambda_i=1\Bigr\}\,.\]

\begin{fact}[Carath\'{e}odory theorem~\cite{Caratheodory}]\label{fact:caratheodory}
Let $X$ be a finite set of points in $\R^d$, and let $y$ be a point
in the convex hull of $X$. Then there exists a subset $Y \subset X$
of size $|Y| \le d+1$ such that $y$ is in the convex hull of $Y$.
\end{fact}

Any hyperplane $H$ partitions a set $X$ of points into two sets: if
$H=\{x:\ip{a,x}=b\}$ then the sets are $\{x \in X:\ip{a,x} \ge b\}$
and $\{x \in X:\ip{a,x}<b\}$. We need the following bound on the maximal
number of ways
a set can be partitioned by hyperplanes.\footnote{A quick way to prove a slightly weaker estimate
is as follows: the $VC$-dimension~\cite{VC71} of
halfspaces is $d+1$. Hence by the VC-dimension
theorem~\cite{VC71,Sauer72,Shelah72} the number of partitions
is at most $\sum_{i=0}^{d+1} {N \choose i} \le N^{d+1}$.}
\begin{fact}[Harding~\cite{Harding67}]
\label{fact:partition_hyperplane}
Let $X$ be a set of $N$ points in $\R^d$. The number
of different ways to partition $X$ into two sets by a hyperplane is
at most $\sum_{i=0}^d {{N-1} \choose i}  \leq N^{d}$.
\end{fact}

\paragraph{Representation theory}
Let $G$ be a finite group. A representation of $G$ (over $\C$) is a
homomorphism $R:G \to \GL(d,\C)$. That is, $R(g)$ for $g \in G$ is a
$d \times d$ nonsingular
complex matrix, and $R(gh) = R(g) R(h)$ for every $g,h \in
G$. The dimension of the representation $R$ is $d$. Two representations
$R,R'$ of $G$ of dimension $d$ are \emph{equivalent} if there exists an
invertible matrix $A$ such that $R'(g) = A^{-1} R(g) A$ for all $g \in
G$. This is denoted as $R \equiv R'$.

A representation $R$ is \emph{unitary}
if $R(g)$ is unitary for all $g \in G$.
\begin{fact}
Any representation of $G$ is equivalent to a unitary representation.
\end{fact}
We will restrict our attention only to unitary representations in this
paper. We note that if $R,R'$ are unitary and equivalent, then there
exists a unitary matrix $A$ such that $R'(g) = A^{-1} R(g) A$.

Let $G$ be a group which acts on a set $X$, that is, $g:X \to X$ is
a permutation of $X$ for all $g \in G$, and $(gh)(x)=g(h(x))$ for all
$g,h \in G$ and $x \in X$. The associated representation $R_X$ maps each
$g \in G$ to the permutation matrix it induces on the set $X$. That is,
$R_X(g)$ is an $|X| \times |X|$ matrix, indexed by $x,y \in X$, defined
as $(R_X(g))_{x,y}=1$ if $g(x)=y$ and $(R_X(g))_{x,y}=0$ otherwise. Note
that $R_X$ is always a unitary representation.

The sum of two representations $R_1:G \to \GL(d_1,\C)$ and $R_2:G \to
\GL(d_2,\C)$ is a representation $R:G \to \GL(d_1+d_2,\C)$ where $R(g)$ is
defined as a block diagonal matrix with two blocks, given by $R_1(g)$ and
$R_2(g)$. For $e \ge 1$ let $e R := R+\cdots+R$ where the sum is over $e$
copies of $R$. A representation $R$ is \emph{reducible} if it is equivalent
to the sum of two representations. Otherwise, the representation
$R$ is \emph{irreducible}. We summarize a few basic properties of
representations below. For details we refer the reader to any standard
book on representation theory, \eg,~\cite{RepresentationTheory}.

\begin{fact}[Maschke's theorem]\label{fact:maschke}
Any representation $R$ of $G$ is equivalent to a sum of irreducible
representations $R \equiv e_1 R_1+\cdots+e_ tR_t$, where $R_1,\ldots,R_t$
are nonequivalent irreducible representations, and $e_i \ge 1$ is the
multiplicity of $R_i$. We have $\sum e_i \dim(R_i) = \dim(R)$.
\end{fact}

\begin{fact}[Schur's lemma]\label{fact:schur}
Let $R$ be a unitary irreducible representation
of $G$ of dimension $d$. Then for any $i,j,k,\ell \in [d]$,
$$
\frac{1}{|G|} \sum_{g \in G} R(g)_{i,j}
\overline{R(g)_{k,\ell}} = \frac{1}{d} \delta_{i,k} \delta_{j,\ell}\,.
$$
Let $R',R''$ be two non-equivalent unitary irreducible
representations of $G$ of dimensions $d',d''$.
Then for any $i,j \in [d']$ and $k,\ell \in [d'']$,
$$
\frac{1}{|G|} \sum_{g \in G} R'(g)_{i,j} \overline{R''(g)_{k,\ell}} = 0\,.
$$
\end{fact}

The trivial representation is given by $\one(g)=1$ for all $g \in G$. An
immediate corollary of Schur's lemma is that
for every representation $R$ which is irreducible
and nontrivial, we have $\sum_{g \in G} R(g)=0$.

The group algebra $\C[G]$ is the linear space of functions $\mu:G \to
\C$. It is often written as $\mu = \sum_{g \in G} \mu(g) \cdot g$. Note
that the distributions over $G$ form
a subset of $\C[G]$. For $\mu \in \C[G]$
and a representation $R$ of $G$, let
$R(\mu) := \sum_{g \in G} \mu(g) R(g)$. If $\mu$ is a distribution,
this is equivalent to $R(\mu) = \E_{g \sim \mu}[R(g)]$. Note that in this
notation, the statistical distance between two distributions $\mu,\mu'$ on
$G$ is $\sd(\mu,\mu')=\|\mu-\mu'\|_1$.


\setcounter{section}{3}\setcounter{theorem}{0}
We first rephrase the conditions for a distribution to be
$X$-uniform, or almost $X$-uniform, in terms of representations. Let
$R_X$ be the representation of the action of $G$ on $X$,
\ie, $R_X(g)_{x,y}=1_{g(x)=y}$. Let $U_G$ denote the uniform distribution
over $G$.

\begin{proposition}\label{prop:rephrase_conditions_dist_Rx}
Let $\mu$ be a distribution on $G$. Then
\begin{enumerate}
\item $\mu$ is $X$-uniform iff $R_X(\mu) = R_X(U_G)$;
\item $\mu$ is almost $X$-uniform with
error $\eps$ iff $\|R_X(\mu)-R_X(U_G)\|_{\infty} \le \eps$.
\end{enumerate}
\end{proposition}

\begin{proof}
The claim is immediate from the definitions of $X$-uniform
and almost $X$-uniform distributions,
since $R_X(\mu)_{x,y}=\Pr_{g \sim \mu}[g(x)=y]$
and $R_X(U_G)_{x,y} = \Pr_{g \in G}[g(x)=y]$.
\end{proof}

The first step is to decompose $R_X$ into its irreducible
representations. Let $R_X \equiv e_0 \one + e_1 R_1 + \cdots + e_t R_t$,
where $R_1,\ldots,R_t$ are unitary nonequivalent non-trivial irreducible
representations, and $e_i$ is the multiplicity of $R_i$ in $R_X$. We next
transform the conditions of \expref{Proposition}{prop:rephrase_conditions_dist_Rx}
to the basis of the irreducible representations.

\begin{proposition}\label{prop:rephrase_conditions_dist_irr}
Let $\mu$ be a distribution on $G$. Then
\begin{enumerate}
\item $\mu$ is $X$-uniform iff $R_i(\mu) = 0$ for all $i \in [t]$;
\item if $\mu$ is almost $X$-uniform with error
$\eps$ then $\|R_i(\mu)\|_{\infty} \le \eps |X|$ for all $i \in [t]$.
\end{enumerate}
\end{proposition}

\begin{proof}
As $\mu$ is a distribution, then $\one(\mu) =
\sum_{g \in G}\mu(g)=1$, and the same holds for $U_G$. Hence always
$\one(\mu)=\one(U_G)$. Thus, $R_X(\mu)=R_X(U_G)$ iff $R_i(\mu)=R_i(U_G)$
for all $i \in [t]$. The first item follows since $R_i(U_G)=0$ for all
$i \in [t]$. To see that, note that by Schur's lemma $$ R_i(U_G)_{j,k}
= \frac{1}{|G|} \sum_{g \in G} R_i(g)_{j,k} = \frac{1}{|G|} \sum_{g
\in G} R_i(g)_{j,k} \overline{\one(g)} = 0 $$ since $R_i$ and $\one$
are nonequivalent unitary irreducible representations. To prove the
second item, let $\mu$ be an almost $X$-uniform distribution with
error $\eps$. By \expref{Proposition}{prop:rephrase_conditions_dist_Rx} this is
equivalent to $\|R_X(\mu)-R_X(U_G)\|_{\infty} \le \eps$. The $L_{\infty}$
norm is not convenient for the basis change required to switch to the
basis of the irreducible representations. We thus switch to the Frobenius
norm, which is trivially bounded by $$\|R_X(\mu)-R_X(U_G)\|_{\Fr}
\le \eps |X|\,.$$ Note that the Frobenius norm is invariant under
unitary change of basis, and as both $R_X$ and $R_1,\ldots,R_t$
are unitary, the basis change can also be assumed to be unitary. We
thus have $$ \sqrt{\sum_{i=1}^t e_i \|R_i(\mu)-R_i(U_G)\|^2_{\Fr}} =
\|R_X(\mu)-R_X(U_G)\|_{\Fr} \le \eps |X|\,,$$ which combined with the
fact that $R_i(U_G)=0$ implies that $\|R_i(\mu)\|_{\infty} \le \eps |X|$.
\end{proof}
\setcounter{section}{4}\setcounter{theorem}{0}
Recall that a distribution $\mu$ is $X$-uniform if $\Pr_{g \sim
\mu}[g(x)=y]=\Pr_{g \in G}[g(x)=y]$ for all $x,y \in X$. We say a set $S$
\emph{supports $X$-uniformity} if there exists a distribution supported on
$S$ which is $X$-uniform. We first establish that this is a purely geometric
property of $S$.

Let $R_X$ be the representation of the action of $G$ on $X$, that
is, \[
R_X(g)_{x,y}
= 1_{g(x)=y}\,.
\]
Let $U=R_X(U_G) = \E_{g \in G} [R_X(g)]$ denote the matrix
which corresponds to the action on $X$ of the uniform distribution over
$G$. We consider these matrices as points in $\R^{d}$ for $d=|X|^2$.

\begin{proposition}\label{prop:uniform-by-convex-hull}
A set $S \subset G$ supports $X$-uniformity iff the convex
hull of the matrices $\{R_X(g): g \in S\}$ contains the matrix $U$.
\end{proposition}
\begin{proof}
A point in the convex hull is given by $M = \sum_{g \in S} \mu(g) \cdot
R_X(g)$ where $\mu(g) \ge 0$ and $\sum_{g \in S} \mu(g)=1$. Thus, each
point in the convex hull corresponds to a distribution $\mu$ over $S$,
and vice versa. Note that $M_{x,y}=\Pr_{g \sim \mu}[g(x)=y]$, hence an
$X$-uniform distribution corresponds to the matrix $U$.
\end{proof}
\setcounter{section}{4}
\section{\texorpdfstring{Almost $X$-uniform distributions support $X$-uniform distributions}{Almost X-uniform distributions support X-uniform distributions}}
\label{sec:almost_support_uniform}

We prove in this section \expref{Theorem}{thm:almost_support_uniform},
which states that
if $\mu$ is an almost $X \times X$-uniform distribution with
small enough error, then there exists an $X$-uniform distribution $\mu'$
supported on the support of $\mu$. We recall the statement of the theorem.

\restate{\expref{Theorem}{thm:almost_support_uniform}}{
Let $\mu$ be a distribution supported on a set $S \subset G$ which is
almost $(X \times X)$-uniform with error $\eps < 0.5 |X|^{-7}$. Then
there exists a distribution $\mu'$ supported on $S$ which is $X$-uniform.
}

Fix such a distribution $\mu$, and let $S$ denote its support, $S=\{g:
\mu(g)>0\}$. Let $R_X$ be the representation of $G$ acting on $X$. By
\expref{Proposition}{prop:uniform-by-convex-hull}, $S$ supports an $X$-uniform
distribution iff $R_X(U_G)=\E_{g \in G}[R_X(g)]$ lies in the convex
hull of $\{R_X(g): g \in S\}$. Assume this is not the case; then there
exists 
%
a hyperplane 
$H$ which passes through $R_X(U_G)$ and such that
all $\{R_X(g): g \in S\}$ lie on one side of $H$.

We first project $H$ into 
%
a hyperplane 
with a simpler representation. Let
$R_X \equiv e_0 \one + e_1 R_1 + \cdots + e_t R_t$ denote the
decomposition of $R_X$ into unitary nonequivalent irreducible
representation, and let $d_i = \dim(R_i)$ denote the dimension of each irreducible representation.
Essentially, we will project $H$ to ``use'' only one copy from each nontrivial
irreducible representation. That is, we will show that $H$ can be
projected to a hyperplane separating $0$ from $\{R_1(g) \times \cdots
\times R_t(g): g \in S\}$.

\begin{proposition}\label{prop:separating_zero_R}
There exists a map $L:G \to \R$ given by
$$
L(g) := \sum_{i \in [t]} \sum_{j,k \in [d_i]}
\lambda_{i,j,k} \cdot R_i(g)_{j,k}
$$
for some coefficients $\{\lambda_{i,j,k}
\in \C: i \in [t], j,k \in [d_i]\}$ such that
\begin{enumerate}
\item $\E_{g \in G}[L(g)]=0$;
\item for all $g \in S$, $L(g)>0$.
\end{enumerate}
\end{proposition}

\begin{proof}
The existence of a hyperplane $H$ separating $R_X(U_G)$ from
$\{R_X(g): g \in S\}$ implies that there exists a map $L':G \to \R$ defined as
$L'(g) = \sum_{x,y \in X} \alpha_{x,y} R_X(g)_{x,y}$
for some real coefficients $\{\alpha_{x,y} \in \R: x,y \in X\}$ such that
$$
L'(g) > \E_{h \in G}[L'(h)]
$$
for all $g \in S$. Applying the linear transformation mapping $R_X$
into the basis of irreducible representations, we get that $L'(g)$
can be expressed as
$$
L'(g) = \sum_{\ell \in [e_0]} \beta_{0,\ell} \one(g)
+ \sum_{i \in [t]} \sum_{j,k \in [d_i]} \sum_{\ell \in [e_i]}
\beta_{i,j,k,\ell} R_i(g)_{j,k}\,,
$$
where $\beta_{0,\ell},\beta_{i,j,k,\ell} \in \C$ are obtained by a linear
transformation (over $\C$) of $\alpha_{x,y}$. Observe that $\E_{g \in
G}[L'(g)] = \sum_{\ell \in [e_0]} \beta_{0,\ell}$ by Schur's lemma,
and define $L(g) := L'(g) - \E[L'(g)]$. Note that $L:G \to \R$ is real
since $L':G \to \R$ was real, that $\E[L]=0$ and that $L(g)>0$ for all
$g \in S$. The coefficient $\lambda_{i,j,k}$ is given by the sum of all
$\beta_{i,j,k,\ell}$ over $\ell \in [e_i]$.
\end{proof}

We may assume without loss of generality that $\E_{g \in G}[L^2(g)]=1$ by multiplying all
coefficients $\lambda_{i,j,k}$ by an appropriate factor. The main idea
is to show that if $\mu$ is almost $X \times X$ uniform, then $\E_{g
\sim \mu}[L^2(g)] \approx \E_{g \in G}[L^2(g)]=1$ while $\E_{g \sim
\mu}[L(g)] \approx \E_{g \in G}[L(g)]=0$. Combining this with a bound
on $\|L\|_{\infty}$ a simple calculation shows that it cannot be the
case that $L(g)>0$ for all $g$ in the support of $\mu$.

The first step is to show that the
coefficients $\lambda_{i,j,k}$ cannot be very large.

\begin{proposition}\label{prop:bound_coefficients_lambda}\hspace{1mm}
\[
\sum_{i \in [t]}\sum_{j,k \in [d_i]} \frac{|\lambda_{i,j,k}|^2}{d_i} = 1\,.
\]
In particular, $|\lambda_{i,j,k}| \le |X|^{1/2}$ for all $i,j,k$.
\end{proposition}

\begin{proof}
We assumed $\E[L^2]=1$, which, since $L$ is
real, implies $\E[|L|^2]=1$. Using Schur's lemma we get
\begin{eqnarray*}
1 &=& \E_{g \in G}\bigl[L(g) \cdot \overline{L(g)}\bigr] \\
&=& \sum_{i,i' \in [t]} \sum_{j,k \in [d_i]} \sum_{j',k' \in [d_{i'}]}
\lambda_{i,j,k} \overline{\lambda_{i',j',k'}} \E_{g \in G}\bigl[R_i(g)_{j,k}
\overline{R_{i'}(g)}_{j',k'}\bigr]\\
& =& \sum_{i \in [t]} \sum_{j,k \in [d_i]} \frac{|\lambda_{i,j,k}|^2}{d_i}\,,
\end{eqnarray*}
and in particular $|\lambda_{i,j,k}|^2 \le d_i \le |X|$.
\end{proof}

An immediate corollary is that $L(g)$ can never be very large.

\begin{corollary}\label{cor:L_bound_infty}
$|L(g)| \le |X|^{2.5}$ for all $g \in G$.
\end{corollary}

\begin{proof}
We have $|R_i(g)_{j,k}| \le 1$ since $R_i$ is unitary, hence $|L(g)|
\le \sum_{i \in [t]} \sum_{j,k \in [d_i]} |\lambda_{i,j,k}|
\le |X|^{2.5}$ since $\sum d_i^2 \le |X|^2$.
\end{proof}

The bound on $|\lambda_{i,j,k}|$ together with the assumption
that $\mu$ is almost $X \times X$-uniform, implies that the first
and second moments of $L$ are approximately
the same under $\mu$ and under the uniform distribution over $G$.

\begin{proposition}\label{prop:L_moments_approx}
Let $\mu$ be a distribution which is
almost $X \times X$-uniform with error $\eps$. Then
\begin{enumerate}
\item $|\E_{g \sim \mu}[L(g)]| \le \eps |X|^{4.5}$;
\item $|\E_{g \sim \mu}[L^2(g)]-1| \le \eps |X|^7$.
\end{enumerate}
\end{proposition}

\begin{proof}
We have
$$
|\E_{g \sim \mu} [L(g)]| \le \sum_{i \in [t]} \sum_{j,k \in [d_i]}
|\lambda_{i,j,k}| |\E_{g \sim \mu} [R_i(g)_{j,k}]|\,.
$$
The bound on the first moment follows since $\sum d_i^2
\le |X|^2$; since $|\lambda_{i,j,k}| \le |X|^{1/2}$ by
\expref{Proposition}{prop:bound_coefficients_lambda}; and since $\mu$
is in particular $X$-uniform with error $\eps |X|$, we have by
\expref{Proposition}{prop:rephrase_conditions_dist_irr} that $|\E_{g \sim
\mu}[R_i(g)_{j,k}]| \le \eps |X|^2$. The bound on the second moment is
proved in a similar way. Recall that we proved the identity $\sum |\lambda_{i,j,k}|^2/d_i=1$ in
\expref{Proposition}{prop:bound_coefficients_lambda}. So
\begin{eqnarray*}
\E_{g \sim \mu} [|L(g)|^2-1] &=& \sum_{i,j,k}
|\lambda_{i,j,k}|^2 \cdot \left(\E_{g \sim \mu}[|R_i(g)_{j,k}|^2] - 1/d_i\right)\\
&+& \sum_{(i,j,k) \ne (i',j',k')} \lambda_{i,j,k}
\overline{\lambda_{i',j',k'}} \cdot
\E_{g \sim \mu}[R_i(g)_{j,k} \overline{R_{i'}(g)_{j',k'}}]\,.
\end{eqnarray*}
To conclude the proof we need to show that $\E_{g \sim
\mu}[|R_i(g)_{j,k}|^2] \approx 1/d_i$ and that $\E_{g \sim
\mu}[R_i(g)_{j,k} R_{i'}(g)_{j',k'}] \approx 0$.

The condition that $\mu$ is
almost $X \times X$-uniform with error $\eps$ is equivalent to
$$
\|R_{X \times X}(\mu) - R_{X \times X}(U_G)\|_{\infty} \le \eps\,.
$$
Switching to the Frobenius norm, this implies
$$
\|R_{X \times X}(\mu) - R_{X \times X}(U_G)\|_{\Fr} \le \eps |X|^2\,.
$$
We now decompose $R_{X \times X}$ to simpler representations, coming
from the irreducible representations of $R_X$. We have that $R_{X
\times X}=R_X \otimes R_X$, which since $R_X$ is real, also gives $R_{X
\times X} = R_X \otimes \overline{R_X}$.  Now, if $R_X \equiv e_0 \one
+ \sum_{i=1}^t e_i R_i$ is the decomposition of $R_X$ into irreducible
unitary nonequivalent representations, we have $$ R_{X \times X} \equiv
e_0^2 \one + \sum_{i=1}^t e_0 e_i (R_i+\overline{R_i}) + \sum_{i,i'=1}^{t}
e_i e_{i'} (R_{i} \otimes \overline{R_{i'}})\,.$$ Note that this is not
the decomposition of $R_{X \times X}$ into irreducible representations,
since $R_i \otimes \overline{R_{i'}}$ is reducible! Nevertheless, we
observe that as the basis change for $R_X$ was unitary, so is the basis
change for $R_{X \times X}$ (since the tensor product of two unitary
matrices is again unitary). In particular, we get that 
\[
\|(R_i \otimes
\overline{R_{i'}})(\mu) - (R_i \otimes \overline{R_{i'}})(U_G)\|_{\Fr}
\le \eps |X|^2\,,
\]
which implies that
$$
\|(R_i \otimes \overline{R_{i'}})(\mu) -
(R_i \otimes \overline{R_{i'}})(U_G)\|_{\infty} \le \eps |X|^2\,.
$$
The matrix $R_i \otimes \overline{R_{i'}}$ is
indexed by $((j,j'),(k,k'))$, where $(R_i \otimes
\overline{R_{i'}})(g)_{(j,j'),(k,k')} = R_i(g)_{j,k}
\overline{R_{i'}(g)_{j',k'}}$. We thus get that for any $i,j,k,i',j',k'$
we have
$$
\left| \E_{g \sim \mu}[R_i(g)_{j,k} \overline{R_{i'}(g)_{j',k'}}]
- \E_{g \in G}[R_i(g)_{j,k} \overline{R_{i'}(g)_{j',k'}}]
\right| \le \eps |X|^2\,.
$$
To conclude, by Schur's lemma
$$
\E_{g \in G}[R_i(g)_{j,k} \overline{R_{i'}(g)_{j',k'}}]
= \frac{1}{d_i} 1_{(i,j,k)=(i',j',k')}\,.
$$
The bound for the second moment now follows by
elementary calculations analogous to the ones for the first moment.
\end{proof}

We are now ready to prove \expref{Theorem}{thm:almost_support_uniform}.

\begin{proof}[Proof of \expref{Theorem}{thm:almost_support_uniform}]
Let $\mu$ be almost $X \times X$ uniform
with error $\eps \le 0.5 |X|^{-7}$.
Summarizing \expref{Corollary}{cor:L_bound_infty} and
\expref{Proposition}{prop:L_moments_approx}, we have
\begin{enumerate}
\item $\|L\|_{\infty} \le |X|^{2.5}$;
\item $E_{g \sim \mu}[L(g)] \le \eps |X|^{4.5}$;
\item $E_{g \sim \mu}[L(g)^2] \ge 1 - \eps |X|^{7}$.
\end{enumerate}
However, since we assumed by contradiction
that $L(g)>0$ for all $g$ in the support of $\mu$, we have
$$
\E_{g \sim \mu}[L(g)^2] \le \|L(g)\|_{\infty}
\cdot \E_{g \sim \mu}[L(g)] \le |X|^{2.5} \cdot \eps |X|^{4.5},
$$
\ie, we have
$$
1 - \eps |X|^7 \le \eps |X|^7,
$$
which is false whenever $\eps < 0.5 |X|^{-7}$.
\end{proof}
\clearpage
\begingroup\raggedright
%
%
%
%
%
\providecommand{\bibhead}[1]{}
%
\expandafter\ifx\csname pdfbookmark\endcsname\relax%
  \providecommand{\tocrefpdfbookmark}{}
\else\providecommand{\tocrefpdfbookmark}{%
   \hypertarget{tocreferences}{}%
   \pdfbookmark[1]{References}{tocreferences}}%
\fi

\tocrefpdfbookmark
\begin{thebibliography}{10}

\bibitem{AA07}\bibhead{AA07}
{\sc Nir Ailon and Noga Alon}: Hardness of fully dense problems.
\newblock {\em Inform. and Comput.}, 205(8):1117--1129, 2007.
\newblock [\epfmtdoi{10.1016/j.ic.2007.02.006}]

\bibitem{AlagicRussell11}\bibhead{AlagicRussell11}
{\sc Gorjan Alagic and Alexander Russell}: Spectral concentration of positive
  functions on compact groups.
\newblock {\em J. Fourier Anal. Appl.}, 17(3):355--373, 2011.
\newblock [\epfmtdoi{10.1007/s00041-011-9174-5}]

\bibitem{AAKMRX07}\bibhead{AAKMRX07}
{\sc Noga Alon, Alexandr Andoni, Tali Kaufman, Kevin Matulef, Ronitt Rubinfeld,
  and Ning Xie}: Testing $k$-wise and almost $k$-wise independence.
\newblock In {\em Proc. 39th STOC}, pp. 496--505. ACM Press, 2007.
\newblock [\epfmtdoi{10.1145/1250790.1250863}]

\bibitem{AlonBabaiItai86}\bibhead{AlonBabaiItai86}
{\sc Noga Alon, L\'{a}szl\'{o} Babai, and Alon Itai}: A fast and simple
  randomized parallel algorithm for the maximal independent set problem.
\newblock {\em J. Algorithms}, 7(4):567--583, 1986.
\newblock [\epfmtdoi{10.1016/0196-6774(86)90019-2}]

\bibitem{AlonGoldreichMansour03}\bibhead{AlonGoldreichMansour03}
{\sc Noga Alon, Oded Goldreich, and Yishay Mansour}: Almost $k$-wise
  independence versus $k$-wise independence.
\newblock {\em Inform. Process. Lett.}, 88(3):107--110, 2003.
\newblock [\epfmtdoi{10.1016/S0020-0190(03)00359-4}]

\bibitem{AustrinHastad11}\bibhead{AustrinHastad11}
{\sc Per Austrin and Johan H{\aa}stad}: Randomly supported independence and
  resistance.
\newblock {\em SIAM J. Comput.}, 40(1):1--27, 2011.
\newblock Preliminary version in
  \href{http://dx.doi.org/10.1145/1536414.1536481}{STOC'09}.
\newblock [\epfmtdoi{10.1137/100783534}]

\bibitem{Cameron95}\bibhead{Cameron95}
{\sc Peter~J. Cameron}: Permutation groups.
\newblock In {\em Handbook of Combinatorics, {V}ol.\ 1}, pp. 611--645.
  Elsevier, Amsterdam, 1995.

\bibitem{Caratheodory}\bibhead{Caratheodory}
{\sc Constantin Carath\'{e}odory}: \"{U}ber den {V}ariabilit\"{a}tsbereich der
  {F}ourier'schen {K}onstanten von positiven harmonischen {F}unktionen.
\newblock {\em Rendiconti del Circolo Matematico di Palermo}, 32(1):193--217,
  1911.
\newblock [\epfmtdoi{10.1007/BF03014795}]

\bibitem{FinucanePeledYaari2012}\bibhead{FinucanePeledYaari2012}
{\sc Hilary Finucane, Ron Peled, and Yariv Yaari}: {A recursive construction of
  $t$-wise uniform permutations}.
\newblock Technical report, 2012.
\newblock \emph{Random Structures and Algorithms}, to appear.
\newblock [\epfmt{arxiv}{1201.4960}]

\bibitem{RepresentationTheory}\bibhead{RepresentationTheory}
{\sc William Fulton and Joe Harris}: {\em Representation Theory: A First
  Course}.
\newblock Volume 129 of {\em Graduate Texts in Mathematics}.
\newblock Springer, 1st edition, 1991.

\bibitem{Harding67}\bibhead{Harding67}
{\sc Edward~Frank Harding}: The number of partitions of a set of {$N$} points
  in {$k$} dimensions induced by hyperplanes.
\newblock {\em Proc. Edinburgh Math. Soc. (2)}, 15(4):285--289, 1967.
\newblock [\epfmtdoi{10.1017/S0013091500011925}]

\bibitem{KaplanNaorReingold05}\bibhead{KaplanNaorReingold05}
{\sc Eyal Kaplan, Moni Naor, and Omer Reingold}: Derandomized constructions of
  $k$-wise (almost) independent permutations.
\newblock {\em Algorithmica}, 55(1):113--133, 2009.
\newblock Preliminary version in
  \href{http://dx.doi.org/10.1007/11538462_30}{RANDOM'05}.
\newblock [\epfmtdoi{10.1007/s00453-008-9267-y}]

\bibitem{KarpPapadimitriou82}\bibhead{KarpPapadimitriou82}
{\sc Richard~M. Karp and Christos~H. Papadimitriou}: On linear
  characterizations of combinatorial optimization problems.
\newblock {\em SIAM J. Comput.}, 11(4):620--632, 1982.
\newblock Preliminary version in
  \href{http://dx.doi.org/10.1109/SFCS.1980.29}{FOCS'80}.
\newblock [\epfmtdoi{10.1137/0211053}]

\bibitem{Kassabov07}\bibhead{Kassabov07}
{\sc Martin Kassabov}: Symmetric groups and expander graphs.
\newblock {\em {Inventiones Mathematicae}}, 170(2):327--354, 2007.
\newblock [\epfmtdoi{10.1007/s00222-007-0065-y}]

\bibitem{KollerMegiddo94}\bibhead{KollerMegiddo94}
{\sc Daphne Koller and Nimrod Megiddo}: Constructing small sample spaces
  satisfying given constraints.
\newblock {\em SIAM J. Discrete Math.}, 7(2):260--274, 1994.
\newblock Preliminary version in
  \href{http://dx.doi.org/10.1145/167088.167168}{STOC'93}.
\newblock [\epfmtdoi{10.1137/S0895480192228140}]

\bibitem{KuehOlsonRockmoreTan01}\bibhead{KuehOlsonRockmoreTan01}
{\sc Ka-Lam Kueh, Timothy Olson, Daniel Rockmore, and Ki-Seng Tan}: Nonlinear
  approximation theory on compact groups.
\newblock {\em J. Fourier Anal. Appl.}, 7(3):257--281, 2001.
\newblock [\epfmtdoi{10.1007/BF02511813}]

\bibitem{KuperbergLovettPeled12}\bibhead{KuperbergLovettPeled12}
{\sc Greg Kuperberg, Shachar Lovett, and Ron Peled}: Probabilistic existence of
  rigid combinatorial structures.
\newblock In {\em Proc. 44th STOC}, pp. 1091--1106. ACM Press, 2012.
\newblock [\epfmtdoi{10.1145/2213977.2214075}]

\bibitem{Maslen98}\bibhead{Maslen98}
{\sc David Maslen}: Efficient computation of {F}ourier transforms on compact
  groups.
\newblock {\em J. Fourier Anal. Appl.}, 4(1):19--52, 1998.
\newblock [\epfmtdoi{10.1007/BF02475926}]

\bibitem{RoyScott09}\bibhead{RoyScott09}
{\sc Aidan Roy and Andrew~J. Scott}: Unitary designs and codes.
\newblock {\em Designs, Codes and Cryptography}, 53(1):13--31, 2009.
\newblock [\epfmtdoi{10.1007/s10623-009-9290-2}]

\bibitem{RubinfeldXie10}\bibhead{RubinfeldXie10}
{\sc Ronitt Rubinfeld and Ning Xie}: Robust characterizations of $k$-wise
  independence over product spaces and related testing results.
\newblock {\em Random Structures \& Algorithms}, 2012 (online).
\newblock Preliminary version in
  \href{http://dx.doi.org/10.1007/978-3-642-14165-2\_48}{ICALP'10}.
\newblock [\epfmtdoi{10.1002/rsa.20423}]

\bibitem{crypto1}\bibhead{crypto1}
{\sc Alexander Russell and Hong Wang}: How to fool an unbounded adversary with
  a short key.
\newblock {\em IEEE Trans. Inform. Theory}, 52(3):1130--1140, 2006.
\newblock Preliminary version in
  \href{http://dx.doi.org/10.1007/3-540-46035-7_9}{EUROCRYPT'02}.
\newblock [\epfmtdoi{10.1109/TIT.2005.864438}]

\bibitem{Sauer72}\bibhead{Sauer72}
{\sc Norbert Sauer}: On the density of families of sets.
\newblock {\em J. Combin. Theory Ser. A}, 13(1):145--147, 1972.
\newblock [\epfmtdoi{10.1016/0097-3165(72)90019-2}]

\bibitem{Shelah72}\bibhead{Shelah72}
{\sc Saharon Shelah}: A combinatorial problem; stability and order for models
  and theories in infinitary languages.
\newblock {\em Pacific J. Math.}, 41(1):247--261, 1972.
\newblock \href{http://projecteuclid.org/euclid.pjm/1102968432}{Project
  Euclid}.

\bibitem{VC71}\bibhead{VC71}
{\sc Vladimir~N. Vapnik and Alexey~Ya. Chervonenkis}: {On the uniform
  convergence of relative frequencies of events to their probabilities}.
\newblock {\em Theory Probab. Appl.}, 16(2):264--280, 1971.
\newblock [\epfmtdoi{10.1137/1116025}]

\bibitem{crypto2}\bibhead{crypto2}
{\sc Serge Vaudenay}: Provable security for block ciphers by decorrelation.
\newblock In {\em Proc. 15th Symp. Theoretical Aspects of Comp. Sci.
  (STACS'98)}, pp. 249--275. Springer, 1998.
\newblock [\epfmtdoi{10.1007/BFb0028566}]

\bibitem{crypto3}\bibhead{crypto3}
{\sc Serge Vaudenay}: Adaptive-attack norm for decorrelation and
  super-pseudorandomness.
\newblock In {\em Proc. 6th Ann. Internat. Workshop on Selected Areas in
  Cryptography (SAC'99)}, pp. 49--61. Springer, 1999.
\newblock [\epfmtdoi{10.1007/3-540-46513-8\_4}]

\bibitem{crypto4}\bibhead{crypto4}
{\sc Serge Vaudenay}: Decorrelation: A theory for block cipher security.
\newblock {\em J. Cryptology}, 16(4):249--286, 2003.
\newblock [\epfmtdoi{10.1007/s00145-003-0220-6}]

\end{thebibliography}
\endgroup
\end{document}
